\section*{第 \ref{ch:g11:stat} 章——描述统计}

\begin{solution}{exo:g11:stat:1}
$n = 8$，和为 $71$，\omterm{def:g11:stat:mean}{平均数} $\bar x = \frac{71}{8} = 8.875$。\omterm{def:g11:stat:median}{中位数}为
第 $4$ 与第 $5$ 个数值的\omterm{prop:g10:coordgeom:midpoint}{中点}，$\frac{8 + 9}{2} = 8.5$。\omterm{def:g11:stat:quartiles}{四分位数}：
$8$ 的四分之一为 $2$，故 $Q_1$ 为第 $2$ 个数值，即 $7$；
四分之三为 $6$，故 $Q_3$ 为第 $6$ 个数值，即 $10$。\omterm{def:g11:stat:quartiles}{四分位距}为
$10 - 7 = 3$.
\end{solution}

\begin{solution}{exo:g11:stat:2}
和为 $40$，故 $\bar x = 5$。平方和：
$4 + 16 + 16 + 16 + 25 + 25 + 49 + 81 = 232$。由简便公式可得
\[
V = \frac{232}{8} - 5^2 = 29 - 25 = 4,
\qquad
\sigma = 2 .
\]
\end{solution}

\begin{solution}{exo:g11:stat:3}
\omterm{def:g11:stat:mean}{平均数}：
\[
\bar x = \frac{1 \times 6 + 2 \times 9 + 3 \times 8 + 4 \times 10
+ 5 \times 9 + 6 \times 8}{50}
= \frac{181}{50} = 3.62 .
\]
\omterm{def:g11:stat:median}{中位数}：累计频数依次为 $6, 15, 23, 33, \dots$，因此
排序后的第 $25$ 与第 $26$ 个数据均为 $4$，故\omterm{def:g11:stat:median}{中位数}为 $4$。
\end{solution}

\begin{solution}{exo:g11:stat:4}
$15$ 个分数之和为 $15 \times 11.2 = 168$；加入
新成员后，\omterm{def:g11:stat:mean}{平均数}为 $\frac{168 + 18}{16} = \frac{186}{16} = 11.625$。
\end{solution}

\begin{solution}{exo:g11:stat:5}
由 $\frac{6 + 8 + x + 12 + 14}{5} = 10$ 得 $40 + x = 50$，故 $x = 10$。
数据 $6, 8, 10, 12, 14$ 与\omterm{def:g11:stat:mean}{平均数} $10$ 的偏差依次为
$-4, -2, 0, 2, 4$，其平方和为 $40$：
$V = \frac{40}{5} = 8$，从而 $\sigma = 2\sqrt2 \approx 2.83$。
\end{solution}

\begin{solution}{exo:g11:stat:6}
由\cref{prop:g11:stat:affine}，取 $a = \frac12$，$b = 4$，新\omterm{def:g11:stat:mean}{平均数}为
$\frac{12.4}{2} + 4 = 10.2$；新\omterm{def:g11:stat:variance}{标准差}为
$\frac{2.5}{2} = 1.25$.
\end{solution}

\begin{solution}{exo:g11:stat:7}
合并\omterm{def:g11:stat:mean}{平均数}是\emph{加权}\omterm{def:g11:stat:mean}{平均数}
\[
\frac{20 \times 10 + 30 \times 15}{50} = \frac{650}{50} = 13,
\]
而非 $12.5$。B 组贡献的数据更多，因此合并\omterm{def:g11:stat:mean}{平均数}被拉向
B 组自身的\omterm{def:g11:stat:mean}{平均数}。
\end{solution}

\begin{solution}{exo:g11:stat:8}
两条线的平均重量都接近目标值（$1.002$ 与 $1.001$ 千克），但
第 2 条线的\omterm{def:g11:stat:variance}{标准差}是另一条线的三倍，其袋装重量波动
大得多，从而明显产生更多欠装与超装的袋子。第 2 条线
应先行处理——需要修正的不只是\omterm{def:g11:stat:mean}{平均数}，
而是一致性。
\end{solution}

\begin{solution}{exo:g11:stat:9}
$\bar x = \frac{70}{10} = 7$;
$V = \frac{540}{10} - 7^2 = 54 - 49 = 5$;
$\sigma = \sqrt5 \approx 2.24$.
\end{solution}

\begin{solution}{exo:g11:stat:10}
\emph{1.} 含 $23$ 时：和为 $50$，\omterm{def:g11:stat:mean}{平均数}为 $5$；\omterm{def:g11:stat:median}{中位数}为第 $5$ 与
第 $6$ 个数值的\omterm{prop:g10:coordgeom:midpoint}{中点} $= 3$。不含 $23$ 时：和为 $27$，\omterm{def:g11:stat:mean}{平均数}为 $3$；\omterm{def:g11:stat:median}{中位数}（9 个数值中的
第 5 个）$= 3$。

\emph{2.} 含 $23$（$n = 10$）时：$Q_1$ 为第 $3$ 个数值（$2$），$Q_3$ 为
第 $8$ 个数值（$4$），\omterm{def:g11:stat:quartiles}{四分位距}为 $2$。不含（$n = 9$）时：$Q_1$ 为
第 $3$ 个数值（$2$），$Q_3$ 为第 $7$ 个数值（$4$），仍为 $2$。

\emph{3.} 单个离群值把\omterm{def:g11:stat:mean}{平均数}从 $3$ 拉高到 $5$，但
不改变\omterm{def:g11:stat:median}{中位数}与\omterm{def:g11:stat:quartiles}{四分位距}：这两者是
\emph{稳健}的。若数据集中可能含有异常值，则应报告
\omterm{def:g11:stat:median}{中位数}与\omterm{def:g11:stat:quartiles}{四分位距}，而非\omterm{def:g11:stat:mean}{平均数}与
\omterm{def:g11:stat:variance}{标准差}。
\end{solution}

\begin{solution}{pb:g11:stat:1}
\textbf{1.} \omterm{def:g11:stat:mean}{平均数} $\bar x = \frac{25}{5} = 5$。偏差
平方依次为 $1, 9, 1, 0, 9$，\omterm{def:g11:stat:variance}{方差}为 $\frac{20}{5} = 4$，故
$\sigma = 2$.

\textbf{2.} 平方的\omterm{def:g11:stat:mean}{平均数}：
$\frac{16 + 64 + 36 + 25 + 4}{5} = 29$；减去
$\bar x^2 = 25$，\omterm{def:g11:stat:variance}{方差}为 $4$。结果相同，计算量减半。

\textbf{3.} \omterm{def:g11:stat:mean}{平均数} $= \frac{3 \times 8 + 5 \times 10 + 2 \times
14}{10} = 10.2$；\omterm{def:g11:stat:variance}{方差} $= \frac{3 \times 64 + 5 \times 100 +
2 \times 196}{10} - 10.2^2 = 108.4 - 104.04 = 4.36$;
$\sigma \approx 2.09$.

\textbf{4.} \omterm{def:g11:stat:mean}{平均数}：$1.8 \times 20 + 32 = 68\,^\circ$F；
$\sigma$：$1.8 \times 3 = 5.4\,^\circ$F——平移 $+32$
不会使数据更分散。

\textbf{5.} 加上 $b$：每个偏差 $x_i + b - (\bar x + b)
= x_i - \bar x$ 不变，故\omterm{def:g11:stat:variance}{方差}不变。将数据乘以
$a$：偏差乘以 $a$，其平方乘以
$a^2$，\omterm{def:g11:stat:variance}{方差}乘以 $a^2$，$\sigma$ 乘以
$\sqrt{a^2} = \abs a$.

\textbf{6.} $z_A = \frac{15 - 12}{2} = 1.5$；
$z_B = \frac{13 - 8}{4} = 1.25$。爱丽丝在所在群体中
遥遥领先：较难考试的 $15$ 分胜过较易考试的 $13$ 分。

\textbf{7.} \omterm{pb:g11:stat:1}{标准分数} $z$ 依次为 $-0.5,\ 1.5,\ 0.5,\ 0,\ -1.5$：\omterm{def:g11:stat:mean}{平均数}
$0$，\omterm{def:g11:stat:variance}{标准差} $1$。一般地，标准化时
先减去 $\bar x$（使\omterm{def:g11:stat:mean}{平均数}变为 $0$，$\sigma$ 不变），
再除以 $\sigma$（使新的 $\sigma$ 等于
$1$）：即第 5 题的两次应用。

\textbf{8.} 由 $z = 2$ 得 $100 + 2 \times 15 = 130$。又
$z(76) = \frac{76 - 100}{15} = -1.6$。

\textbf{9.} 该婴儿的体重比
\emph{同龄婴儿}的平均体重低 $2.5$ 个\omterm{def:g11:stat:variance}{标准差}——相对
同龄群体而言轻得惊人。千克数无法跨年龄比较（所有
婴儿都会增重）；\omterm{pb:g11:stat:1}{标准分数}把每个孩子与其所在年龄组对比，
为所有年龄提供同一把尺子。

\textbf{10.} \omterm{def:g11:stat:mean}{平均数} $5$；平方的\omterm{def:g11:stat:mean}{平均数} $\frac{622}{10} = 62.2$，
\omterm{def:g11:stat:variance}{方差} $37.2$，$\sigma \approx 6.10$；
$z(23) = \frac{18}{6.10} \approx 2.95$：略低于报警
阈值——检验几乎毫无反应。

\textbf{11.} 去掉 $23$ 后：\omterm{def:g11:stat:mean}{平均数} $3$，\omterm{def:g11:stat:variance}{方差}
$\frac{93}{9} - 9 \approx 1.33$，$\sigma \approx 1.15$；
以去除离群值后的尺子衡量，$23$ 的\omterm{pb:g11:stat:1}{标准分数}为
$z \approx \frac{20}{1.15} \approx 17$：一个怪物。异常值
同时抬高了\omterm{def:g11:stat:mean}{平均数}与 $\sigma$，使自身得以伪装——
这正是 \cref{exo:g11:stat:10} 所建议的：
数据受污染时应使用稳健统计量（\omterm{def:g11:stat:median}{中位数}、\omterm{def:g11:stat:quartiles}{四分位距}）。

\textbf{12.} 在\omterm{def:g11:stat:variance}{方差}和式
$\sigma^2 = \frac1N \sum (x_i - \bar x)^2$ 中，只保留
满足 $\abs{x_i - \bar x} \geq k\sigma$ 的项：共有 $pN$
个，每项都至少为 $(k\sigma)^2$，于是
$\sigma^2 \geq \frac{1}{N} \cdot pN \cdot k^2\sigma^2 =
p\,k^2\sigma^2$：两边除以 $k^2 \sigma^2$（离散程度非零）：
$p \leq \frac{1}{k^2}$。

\textbf{13.} 超出 $2\sigma$ 的部分，在任何数据集中至多占
$\frac14 = 25\,\%$；超出 $3\sigma$ 的部分至多为
$\frac19 \approx 11\,\%$。第 1 题的数据中，
超出 $2\sigma = 4$（距\omterm{def:g11:stat:mean}{平均数} $5$）即指数值 $\leq 1$ 或 $\geq 9$：
五个数值均不在此范围内——$0\,\% \leq 25\,\%$，与前述结论相符。

\textbf{14.} 落在 $\intcc{49.7}{50.3}$ 之外，即超出
$3\sigma$：至多占产量的 $\frac19$，\emph{无论}
分布的形状如何——切比雪夫不要求其他任何条件，只需
\omterm{def:g11:stat:mean}{平均数}与一个 $\sigma$。钟形分布的产量集中得好得多
（仅约 $0.3\,\%$ 超出 $3\sigma$）：正态分布
将在十二年级学习。

\textbf{15.} $-19\,\%$ 比\omterm{def:g11:stat:mean}{平均数}低 $24$ 个百分点（\omterm{def:g11:stat:mean}{平均数}为 $5$）。
基金 B：$24 = 2\sigma_B$，切比雪夫允许至多
$\frac14$ 的年份如此糟糕。基金 A：$24 = 12\sigma_A$，
至多占 $\frac{1}{144}$ 的年份。平均收益率相同，风险却不可同日而语——
$\sigma$ 即金融领域所说的\emph{风险}（或波动率）。

\textbf{16.} 围绕 $\bar x$ \omterm{def:g10:algebra:expand}{展开}：
$(x_i - c)^2 = \left((x_i - \bar x) + (\bar x - c)\right)^2
= (x_i - \bar x)^2 + 2(\bar x - c)(x_i - \bar x) +
(\bar x - c)^2$。取平均后中间项消去（与\omterm{def:g11:stat:mean}{平均数}的偏差
之和为零），于是剩下
$f(c) = \sigma^2 + (\bar x - c)^2$——仅在
$c = \bar x$ 时取\omterm{def:g10:functions:extrema}{最小值}，\omterm{def:g10:functions:extrema}{最小值}为 $\sigma^2$。\omterm{def:g11:stat:mean}{平均数}是
最小二乘中心，正如\omterm{def:g11:stat:median}{中位数}是最小绝对离差
中心（见\cref{pb:g10:stats:1}）。

\textbf{17.} 直接计算：与 $6$ 的偏差依次为
$-2, 2, 0, -1, -4$，平方依次为 $4, 4, 0, 1, 16$：
$f(6) = \frac{25}{5} = 5$。用公式：
$\sigma^2 + (5 - 6)^2 = 4 + 1 = 5$。两者一致。

\textbf{18.} 两组数据的\omterm{def:g11:stat:mean}{平均数}均为 $7$；\omterm{def:g11:stat:variance}{方差}：
$P$：$\frac{16 + 4 + 0 + 4 + 16}{5} = 8$；
$Q$：$\frac{36 + 0 + 0 + 0 + 36}{5} = 14.4$：不同的离散程度
隐藏于同一\omterm{def:g11:stat:mean}{平均数}之后。对于 $\{7 - a, 7, 7, 7, 7 + a\}$，\omterm{def:g11:stat:variance}{方差}
$\frac{2a^2}{5} = 8$ 可得 $a = \sqrt{20} \approx 4.47$：
数据集 $\{2.53, 7, 7, 7, 11.47\}$ 恰好具有与
$P$ 相同的\omterm{def:g11:stat:mean}{平均数}与\omterm{def:g11:stat:variance}{方差}，而分布形态截然不同——
汇总即是压缩，压缩必然丢失信息。

\textbf{19.} $\sigma^2 = 0$ 表示平方和
$(x_i - \bar x)^2$ 为零；平方 $\geq 0$，故
和为零，则每一项都必为零：$x_i = \bar x$ 对所有 $i$ 成立
——所有数值相等。

\textbf{20.} 中心说明数据位于何处，离散程度说明
数据散布得有多开；\omterm{pb:g11:stat:1}{标准分数}为所有数据集、
考试、年龄与货币提供同一把尺子；切比雪夫将 $\sigma$ 化为
无条件的保证；而第 18 题提醒我们保持审慎
——寥寥几个数字永远无法包含全部信息。下一
年，钟形曲线将把这把尺子磨成可给出精确
\omterm{def:g10:proba:distribution}{概率}的工具。
\end{solution}
